%=====================================================================
\section{Lecture 8: Mori's projection, the generalized Langevin equation and the retarded self-energy}
\label{sec:lec8}

Nakajima--Zwanzig projects \emph{states}. Mori~\cite{Mori1965} projects
\emph{operators}: pick a few ``slow'' Heisenberg operators, and derive an exact equation
of motion for them in which everything else appears as a memory term and a random force.
The output is a generalized Langevin equation (GLE), and its Laplace transform is again a
Dyson equation, now for correlation functions. This lecture derives the GLE, works it out
for the qubit's lowering operator, and establishes two precise relations: Mori on the
\emph{complete} operator basis is the Heisenberg-picture adjoint of NZ, and Mori's memory
function for a single operator is the retarded self-energy of that operator's Green's
function.

\subsection{The generalized Langevin equation}

Work in the Heisenberg picture: $\dot A(t)=i[H,A(t)]\equiv i\mathsf LA(t)$, so $A(t)=\ee^{i\mathsf Lt}A$. Choose
an inner product $(X,Y)$ on operators, a set of operators $\{A_j\}$ (one, for simplicity;
the matrix generalization is immediate) and the projector
\begin{equation}
  \mathsf PX=A\,(A,A)^{-1}(A,X),\qquad \mathsf Q=1-\mathsf P .
\end{equation}
The derivation parallels \cref{sec:nz-derivation} with the Dyson-type identity
$\ee^{i\mathsf Lt}=\ee^{i\mathsf Q\mathsf Lt}+\int_0^t\dd s\,\ee^{i\mathsf L(t-s)}\,i\mathsf P\mathsf L\,\ee^{i\mathsf Q\mathsf Ls}$ applied to
$i\mathsf LA$~\cite[Ch.~8]{ZwanzigBook}, \cite[Secs.~5.2--5.3]{Forster1975}. The result is exact:
\begin{equation}
  \boxed{\;\dot A(t)=i\Omega\,A(t)-\int_0^t\!\dd s\,M(t-s)\,A(s)+F(t),\qquad
  F(t)=\ee^{i\mathsf Q\mathsf Lt}\,\mathsf Q\,i\mathsf LA,\;}
  \label{eq:GLE}
\end{equation}
with the \emph{frequency} $i\Omega=(A,i\mathsf LA)(A,A)^{-1}$ and the \emph{memory function}
\begin{equation}
  M(t)=\big(F(0),F(t)\big)\,(A,A)^{-1}.
  \label{eq:memory-fn}
\end{equation}
The random force lives in the orthogonal complement, $(A,F(t))=0$ for all $t$, which is what
makes the correlation function $C(t)\equiv(A,A(t))$ obey the closed equation
\begin{equation}
  \dot C(t)=i\Omega\,C(t)-\int_0^t\!\dd s\,M(t-s)\,C(s)
  \quad\Longrightarrow\quad
  \hat C(z)=\frac{C(0)}{z-i\Omega+\hat M(z)} .
  \label{eq:MoriDyson}
\end{equation}
\Cref{eq:memory-fn}, the statement that the memory kernel is the autocorrelation of the
random force, is Kubo's \emph{second fluctuation--dissipation theorem}~\cite{Kubo1966}.
Everything depends on the inner product. The two standard choices are the
\emph{Kubo--Mori} product $(X,Y)=\beta^{-1}\int_0^\beta\dd\lambda\avg{X^\dagger\ee^{-\lambda H}Y\ee^{\lambda H}}$, for which
$C(t)$ is the canonical (Kubo) correlation and $\mathsf L$ is self-adjoint, and the
\emph{symmetrized} product $(X,Y)=\tfrac12\avg{\{X^\dagger,Y\}}$, which is simpler for fermions and
spins; at the order we work, the two agree. In both, $\avg\cdot$ is the thermal average in
the \emph{correlated} global equilibrium state, not the factorized state of
\cref{sec:lec5}; this difference will show up.

\subsection{Mori on the complete basis is the adjoint of Nakajima--Zwanzig}

\begin{proposition}\label{prop:mori-nz}
Let $\{E_\mu\}_{\mu=1}^{d^2}$ be an orthonormal basis of system operators,
$\Tr_S(E_\mu^\dagger E_\nu)=\delta_{\mu\nu}$, and take the inner product
$(X,Y)_B\equiv\Tr\big[X^\dagger Y\,(\Id_S\otimes\rhoB)\big]$ on operators of $S\otimes B$. Mori's projector onto
the set $\{E_\mu\otimes\Id_B\}$ is the Hilbert--Schmidt adjoint of the NZ projector
\eqref{eq:P}: $\mathsf P=\Pc^\dagger$, with $\Pc^\dagger Y=\Tr_B[(\Id_S\otimes\rhoB)Y]\otimes\Id_B$. Consequently the
Mori memory matrix for this set is the adjoint of the NZ kernel, $M(t)=-K(t)^\dagger$ (for the
Heisenberg generator $\mathsf L^\dagger$ in place of $\mathsf L$), and the GLE \eqref{eq:GLE} for the
Heisenberg operators $E_\mu(t)$ is the adjoint of the NZ equation.
\end{proposition}
\begin{proof}
$\mathsf PX=\sum_\mu E_\mu\,(E_\mu,X)_B=\sum_\mu E_\mu\Tr[(E_\mu^\dagger\otimes\rhoB)X]
=\sum_\mu E_\mu\Tr_S\!\big[E_\mu^\dagger\Tr_B[(\Id\otimes\rhoB)X]\big]=\Tr_B[(\Id\otimes\rhoB)X]\otimes\Id_B$ by completeness, and
$(Y,\Pc X)_{\rm HS}=\Tr[Y^\dagger(\rhoB\otimes\Tr_BX)]=\Tr_S\big[\Tr_B[(\Id\otimes\rhoB)Y]^\dagger\,\Tr_BX\big]=(\Pc^\dagger Y,X)_{\rm HS}$.
The rest is the general fact that taking the adjoint of \cref{eq:NZ} (with
$\Lc^\dagger=i\mathsf L$ and $\Pc^\dagger=\mathsf P$) gives \cref{eq:GLE} for the full basis, since the
derivations are identical line by line.
\end{proof}

So the two projection techniques are not merely analogous: with the \emph{same}
reference state and the \emph{complete} set of system operators, they are the same
equation written in the two pictures. The novelty of Mori's scheme is in the two freedoms
it adds: a \emph{reduced} set of slow variables (one does not need all $d^2$ of them), and an
inner product that weights with the \emph{interacting} equilibrium state, so that the
reference state already knows about the system--bath correlations. The first is what
makes Mori useful for hydrodynamics and for Green's functions; the second is what makes
its results differ from factorized-initial-state NZ even at second order, as we now see.

\subsection{Example: Mori for the qubit's lowering operator}

\begin{exbox}[Example 8.1: $A=\sigma_-$ and the retarded self-energy]
Take $A=\sigma_-$ with the symmetrized product in the state $\rho_S^{\rm eq}\otimes\rhoB$ (the
lowest-order approximation to the correlated equilibrium state), $\rho_S^{\rm eq}\propto\ee^{-\beta H_S}$. Then
$(\sigma_-,\sigma_-)=\tfrac12\avg{\{\sigma_+,\sigma_-\}}=\tfrac12$. Using $[\sz,\sigma_-]=-2\sigma_-$ and $[\sx,\sigma_-]=-\sz$,
\begin{equation}
  i\mathsf L\sigma_-=i[H,\sigma_-]=-i\omega_0\sigma_--i\big(2\cos\theta\,\sigma_-+\sin\theta\,\sz\big)B .
\end{equation}
The first term is along $A$, so $i\Omega=-i\omega_0$; the second is orthogonal to $A$ (it has
$\avg B=0$) and is the random force, $F(0)=-iX_-B$ with $X_-\equiv2\cos\theta\sigma_-+\sin\theta\sz$. To
second order, $F(t)$ evolves freely: $F(t)=-iX_-(t)B(t)$ with
$X_-(t)=2\cos\theta\sigma_-\ee^{-i\omega_0t}+\sin\theta\sz$. The memory function \eqref{eq:memory-fn} is
$M(t)=2(F(0),F(t))=\avg{\{X_+B,X_-(t)B(t)\}}$ with $X_+=X_-^\dagger$. The system and bath
averages factorize, $\avg{X_+BX_-(t)B(t)}=\avg{X_+X_-(t)}\,C(-t)$ and
$\avg{X_-(t)B(t)X_+B}=\avg{X_-(t)X_+}\,C(t)$; with $\avg{\sigma_+\sigma_-}=p_e$, $\avg{\sigma_-\sigma_+}=p_g$ (excited and
ground populations) and $\avg{\sz^2}=1$,
\begin{equation}
  M(t)=2\sin^2\theta\,S(t)+4\cos^2\theta\,\ee^{-i\omega_0t}\Big[S(t)-\tfrac i2\tanh\tfrac{\beta\omega_0}{2}\,\chi(t)\Big]\qquad(t>0),
  \label{eq:M-qubit}
\end{equation}
where we used $p_gC(t)+p_eC(-t)=S-\tfrac i2(p_g-p_e)\chi$ and $p_g-p_e=\tanh(\beta\omega_0/2)$.

Now the Green's function. With $G^R(t)=-i\theta(t)\avg{\{\sigma_-(t),\sigma_+\}}=-2i\theta(t)C(t)$, the Laplace
transform of \cref{eq:MoriDyson} at $z=-i\omega+0^+$ gives
\begin{equation}
  G^R(\omega)=\frac{1}{\omega-\omega_0-\Sigma^R(\omega)},\qquad
  \boxed{\;\Sigma^R(\omega)=-i\hat M(-i\omega+0^+)=-i\int_0^\infty\!\dd t\,\ee^{i\omega t}M(t).\;}
  \label{eq:SigmaR}
\end{equation}
Mori's memory function \emph{is} the retarded self-energy of the spin-flip propagator.
Evaluating at $\omega=\omega_0$ with \cref{eq:M-qubit}:
\begin{align}
  -\Im\Sigma^R(\omega_0)&=\Re\int_0^\infty\!\ee^{i\omega_0t}M(t)\,\dd t
  =\sin^2\theta\,S(\omega_0)+2\cos^2\theta\,S(0)=\tfrac12\Gamma_1+\Gamma_\varphi=\Gamma_2 ,\\
  \Re\Sigma^R(\omega_0)&=\Im\int_0^\infty\!\ee^{i\omega_0t}M(t)\,\dd t
  =\underbrace{2\sin^2\theta\int_0^\infty\!S(t)\sin\omega_0t\,\dd t}_{\delta\omega\ \text{of \cref{eq:lamb}}}
  \;-\;2\cos^2\theta\tanh\tfrac{\beta\omega_0}{2}\;\chi(\omega{=}0).
\end{align}
The decay rate is exactly the Bloch--Redfield $\Gamma_2$ of \cref{eq:rates-T2}: the
transverse part of the memory sampled at $\omega_0$ gives $\Gamma_1/2$ (the $(\sigma_-\ee^{-i\omega_0t})$ term
has been shifted to zero frequency by the $\ee^{i\omega_0t}$, which is why it yields $S(0)$), and the
longitudinal part gives $\Gamma_\varphi$. The real part contains the Lamb shift $\delta\omega$ of
\cref{sec:lec5} \emph{plus} a term that \cref{sec:lec5} did not have: a static shift
$-2\cos^2\theta\tanh(\beta\omega_0/2)\chi(0)$, with $\chi(0)=\tfrac2\pi\int_0^\infty J(\omega)/\omega\,\dd\omega=2\lambda$ twice the
reorganization energy. It is a Franck--Condon (Stokes) shift: in the \emph{equilibrium}
state the bath is displaced in opposite directions for the two qubit levels, so a
vertical transition from the thermally more populated level costs $\omega_0\mp4\lambda$, and the
thermal average of the two line positions is $\omega_0-4\lambda\tanh(\beta\omega_0/2)$. Mori sees it because
its reference state is correlated; NZ from a factorized initial state sees it only
dynamically, through the $\cos\theta\sin\theta\,\chi$ cross terms in the first column of
\cref{eq:K2-bloch} that the secular approximation discards. Same physics, different
bookkeeping of the initial condition.
\end{exbox}

\begin{keybox}[Key point 8.1]
For a single projected operator $A$ with $(A,A)$ normalized, Mori's memory function is
$-i$ times the retarded self-energy of the Green's function $G^R_{AA^\dagger}$, and the
frequency $\Omega$ is the static (Hartree--Fock-like) part of that self-energy. Iterating
the projection (project $F$ onto itself, etc.) produces Mori's continued fraction, which
is the recursion method for $G^R$~\cite[Ch.~8]{ZwanzigBook}, \cite[Sec.~5.3]{Forster1975};
for the resonant-level model the construction reproduces the exact hybridization
self-energy (see the companion working note).
\end{keybox}

\subsection{What Mori does \emph{not} give you}

\Cref{eq:MoriDyson} determines the retarded sector. In equilibrium the FDT then fixes
the Keldysh component, $G^K=\coth(\beta\omega/2)(G^R-G^A)$ for bosonic operators and
$(1-2f)(G^R-G^A)$ for fermionic ones, so nothing is lost. Out of equilibrium,
occupations, currents and noise live in $G^<$ (equivalently $G^K$), which the retarded
function alone does not determine. Three routes supply the missing information:
\begin{enumerate}[itemsep=1pt]
\item \emph{Density-matrix projection} (NZ/TCL, \cref{sec:lec5,sec:lec6}), which evolves
  the one-time state and therefore carries occupations by construction; by
  \cref{prop:mori-nz} this is Mori on the complete basis with a factorized reference state.
\item \emph{Kadanoff--Baym equations}~\cite{KadanoffBaym,StefanucciVanLeeuwen} for the
  two-time $G^{\lessgtr}(t,t')$, closed by the generalized Kadanoff--Baym ansatz
  (GKBA)~\cite{Lipavsky1986}, $G^{\lessgtr}(t,t')\approx i[G^R(t,t')G^{\lessgtr}(t',t')-G^{\lessgtr}(t,t)G^A(t,t')]$, which
  plays the role of a projection onto one-time quantities.
\item \emph{Time-dependent Mori projections} with nonstationary reference
  states~\cite[Ch.~2]{Grabert1982}, which produce two-time memory kernels and are closer in
  spirit to Kadanoff--Baym.
\end{enumerate}

\subsection{GLE $\leftrightarrow$ Keldysh saddle point}

For a linear bath coupled to a continuous coordinate $x$ (Caldeira--Leggett), Example~7.1(ii)
produced from the influence functional the quantum Langevin equation
$m\ddot x^\cl+V'(x^\cl)-\int^t\chi(t-s)x^\cl(s)\dd s=\xi(t)$, $\avg{\xi\xi}=S$. This is Mori's
\cref{eq:GLE} for $A=x$ (with the momentum as a second variable), with
$M\leftrightarrow\chi$ and $F\leftrightarrow\xi$; Kubo's second FDT \eqref{eq:memory-fn} and the
equilibrium relation $D^K=\coth(\beta\omega/2)(D^R-D^A)$ between the bath's Keldysh and retarded
components are the same statement~\cite[Sec.~3.2]{Kamenev2011}, \cite[Ch.~3]{Gardiner2004}.
For the qubit there is no such classical limit, but the structure is unchanged: the
memory function \eqref{eq:M-qubit} is a combination of $S$ and $\chi$ with the system's
own thermal weight $\tanh(\beta\omega_0/2)$ multiplying the response.

\subsection{Exercises}

\begin{exercise}\label{ex:8-1}
Derive \cref{eq:GLE} from the Dyson identity and prove $(A,F(t))=0$. Show that
$M(t)=(F,F(t))(A,A)^{-1}$ requires $\mathsf L$ to be self-adjoint for the chosen product, and
check this for the Kubo--Mori and symmetrized products in a thermal state.
\end{exercise}

\begin{exercise}\label{ex:8-2}
Complete the derivation of \cref{eq:M-qubit} (the four anticommutator terms) and verify the
rate and shift formulas. Then repeat with $A=\sz-\avg{\sz}$ and show that the memory
function reproduces $\Gamma_1$ (use the FDT) and that the inhomogeneous term is absent (why?).
What goes wrong if one projects on $\sz$ itself?
\end{exercise}

\begin{exercise}\label{ex:8-3}
Prove \cref{prop:mori-nz} in detail for a qubit using the basis
$E_\mu=\sigma_\mu/\sqrt2$, and show that the Mori memory matrix at second order is
$-K^{(2)}(t)^T$ in the coefficient convention (the transpose, because the matrices are
real).
\end{exercise}

\begin{exercise}\label{ex:8-4}
Mori's continued fraction: project the random force of Example~8.1 onto itself and
write the second-level memory function. For the Lorentzian bath, show that the
continued fraction terminates after two levels in the rotating-wave approximation and
reproduces the exact $c_e(t)$ of \cref{sec:tcl-singular}.
\end{exercise}

\clearpage
